% !TeX root = ../main.tex
\subsection{Пункт 7. Исследование характеристик при использовании единого КИХ-дециматора}

В данном разделе представлены результаты выполнения седьмого пункта задания, в котором многоскоростная система фильтрации на базе каскада CIC-фильтра и КИХ-корректора заменяется на единый фильтр с конечной импульсной характеристикой (КИХ-фильтр). Спроектированный КИХ-дециматор обладает требуемыми параметрами: неравномерность в полосе пропускания 0.1 дБ и ослабление в полосе задерживания 60 дБ. Понижение частоты дискретизации ($R=10$) интегрировано непосредственно в процесс фильтрации.

\textbf{Зависимость амплитуды побочных составляющих от разрядности ЦГ}

Для новой архитектуры было повторено исследование влияния эффектов квантования цифрового гетеродина (режим \texttt{nco\_type = 'fixed'}) на чистоту выходного спектра. Диапазон изменения разрядности ЦГ составил от 2 до 11 бит. Полученная зависимость динамического диапазона (SFDR) от количества бит ЦГ представлена на рисунке \ref{fig:p7_sfdr_bits}.

\begin{figure}[H]
    \centering
    \includegraphics[width=0.8\textwidth]{pynkt_7_3.png}
    \caption{Зависимость SFDR от количества бит ЦГ при использовании единого КИХ-дециматора}
    \label{fig:p7_sfdr_bits}
\end{figure}

\textbf{Сравнительный анализ и выводы:}
Полученный график для единого КИХ-дециматора кардинально отличается от аналогичного графика для CIC-фильтра. Здесь наблюдается классическая теоретическая картина: при увеличении разрядности цифрового гетеродина от 3 до 9 бит значение SFDR линейно возрастает с $\approx 23$ дБн до $\approx 83$ дБн (в среднем улучшаясь на 6 дБ с каждым добавленным битом). При $n_g \ge 10$ бит рост прекращается, и SFDR стабилизируется на уровне $\approx 74$--$75$ дБн, упираясь в физический шумовой порог тракта (входной тепловой шум и шумы квантования 9-битного АЦП).

Столь разительное отличие в поведении двух схем фильтрации обусловлено архитектурными особенностями обработки высокочастотных паразитных гармоник и явлением спектрального наложения (aliasing):

\begin{itemize}
    \item \textbf{Проблема CIC-фильтра:} Каскад интеграторов в фильтре Хогенауэра имеет специфическую амплитудно-частотную характеристику с широкими боковыми лепестками и плавными спадами. Мощные высокочастотные паразитные гармоники, порождаемые квантованием гетеродина, недостаточно подавляются интеграторами до момента децимации. В процессе понижения частоты дискретизации эти неподавленные гармоники алиасятся в рабочую полосу частот. Этот интенсивный алиасинг создавал высокий уровень шумового фона, маскируя любые улучшения от повышения разрядности ЦГ.
    \item \textbf{Преимущество единого КИХ-дециматора:} Одиночный КИХ-фильтр осуществляет строгую низкочастотную фильтрацию с высоким гарантированным подавлением (60 дБ) и крутым спадом характеристики на \textit{полной} частоте дискретизации АЦП ($f_s$) еще до отбрасывания отсчетов. Это надежно подавляет широкополосный шум квантования гетеродина и предотвращает его попадание в базовую полосу при децимации. 
\end{itemize}

Таким образом, использование качественного КИХ-фильтра позволяет полностью раскрыть потенциал многоразрядного цифрового гетеродина, обеспечивая значительно более высокий динамический диапазон системы по сравнению с экономичным, но спектрально уступающим CIC-дециматором.